%% ma: L10-B13-0013
%% lop: 10
%% chuong: 5
%% bai: 13
%% dang: 10.13.2
%% loai: TN4
%% muc: TH
%% dap_an: "C"
%% nguon: "chuyen-de-cac-so-dac-trung-cua-mau-so-lieu-khong-ghep-nhom-toan-10-chuong-trinh-moi.pdf (GV. Phan Nhật Linh) (P02), Bài 02 Dạng 1 Phần I Câu 3"
%% kien_thuc: "Lớp 10 Bài 13 (số trung bình)"
%% trang_thai: chinh-thuc
%% ghi_chu: "Tự duyệt theo yêu cầu giáo viên 22/10/2026."
\begin{ex}
Cho bảng tần số về sản lượng chè thu được trong $1$ năm (kg/sào) của $20$ hộ gia đình:
\begin{center}
\begin{tabular}{|c|c|c|c|c|c|c|c|}\hline
Sản lượng & $111$ & $112$ & $113$ & $114$ & $115$ & $116$ & $117$\\\hline
Tần số & $1$ & $3$ & $4$ & $5$ & $4$ & $2$ & $1$\\\hline
\end{tabular}
\end{center}
Số trung bình của mẫu số liệu trên là
\choice{$114$}{$114{,}5$}{\True $113{,}9$}{$113{,}5$}
\loigiai{$\overline x=\dfrac{1\cdot111+3\cdot112+4\cdot113+5\cdot114+4\cdot115+2\cdot116+1\cdot117}{20}=\dfrac{2278}{20}=113{,}9$. Chọn C.}
\end{ex}
